%% Formule de l'operateur : Cor_rgd_assiet
%% Dans ce fichier, les �ventuels accents sont cod�s
%% en ISO-8859-15 (aka latin9)
%% La ligne suivante donne trois exemples 
%% e accent aigu : �, e accent grave : �, e-dans-l'o : �
%Les angles d'attitudes de l'instrument sont calcul�s 
%� partir de ceux de la plate-forme et du d�calage :
%\[ \begin{cases}
%\phi_i &= \phi_p + \Delta\phi \\
%\theta_i &= \theta_p + \Delta\theta
%\end{cases}
%\]
\newcommand{\JDA}{\ensuremath{JDA}}
\newcommand{\THETAZERO}{\ensuremath{\Theta_0}}
\newcommand{\PHIZERO}{\ensuremath{\Phi_0}}
\newcommand{\AMUZERO}{}
\newcommand{\CAP}{\ensuremath{\psi}}
\newcommand{\HZERO}{\ensuremath{h_0}}
Une fonction d'�ph�m�ride donne, � partir du
num�ro du jour dans l'ann�e (argument \JDA), la position du soleil :
\THETAZERO (angle solaire z�nithal), \PHIZERO (angle solaire azimutal),
\AMUZERO
\HZERO (hauteur angulaire du Soleil sur l'horizon, $= 90-\THETAZERO$)


Puis la correction est appliqu�e sous la forme :
\[
SW_c = \frac{SW_d}{Xn Xas}
\]
\newcommand{\OFROL}{\Delta\phi}
\newcommand{\OFTAG}{\Delta\theta}
\newcommand{\DSIR}{\sin(\phi_p)}
\newcommand{\DCOR}{\cos(\phi_p)}
\newcommand{\DSIP}{\sin(\theta_p)}
\newcommand{\DCOP}{\cos(\theta_p)}
\newcommand{\DSIA}{\sin(\PHIZERO - \CAP)}
\newcommand{\DCOA}{\cos(\PHIZERO - \CAP)}
\newcommand{\DTAH}{\tan(\HZERO)}
\newcommand{\DCUN}{-\DCOR\DSIA+\DSIR\DCOA\DSIP}
\newcommand{\DCDEUX}{\DCOP\DCOA}

\noindent o� $Xn$ correspond � la correction d'incidence solaire sur le capteur :
\[
    Xn = \sin(\OFTAG)X_0 + \sin(\OFROL)X_1 + \sqrt{1 - \sin(\OFROL)^2 - \sin(\OFTAG)^2}X_2
\]
avec :\newline
$X_0 = \frac{-(\DCOP\DCOA)}{\DTAH-\DSIP} $\newline
\newline
$X_1 = \frac{-(\DSIP\DSIR\DCOA+\DCOR\DSIA)}{\DTAH+\DCOP\DSIR} $\newline
\newline
$X_2 = \frac{-(\DSIP\DCOR\DCOA-\DSIR\DSIA)}{\DTAH+\DCOP\DCOR} $\newline

et $Xas$ correspond � la correction de cosinus :
\[ Xas = 1 - \frac{3.9 - 12.5 C + 17.5 C^2 - 8.9 C^3}{100} 
\qquad \text{avec } C = \cos(\THETAZERO)
\]

